\section{A sequence defined by a recurrence}\label{ch05:sec:recurrence}
A \emph{sequence} is an endless list of numbers $x_0,x_1,x_2,\ldots$, one for each natural number. Each $x_n$ is a \emph{term} of the sequence, and the subscript $n$ is its position, counting from $0$. A \emph{recurrence} defines a sequence step by step: it gives the first term and a rule that computes each later term from the earlier ones. Let
\[
x_0=0\qquad\text{and}\qquad x_{n+1}=\frac{x_n+1}{2}\quad\text{for every } n\in\NN.
\]
Read the notation carefully: $x_{n+1}$ is the term after $x_n$, whereas $x_n+1$ is the current term plus one. So the rule says: add one to the current term and halve the result. Starting from $x_0=0$,
\[
x_1=\frac{0+1}{2}=\frac12,\qquad x_2=\frac{\frac12+1}{2}=\frac34,\qquad x_3=\frac{\frac34+1}{2}=\frac78.
\]
The terms fall short of $1$ by $1$, $\tfrac12$, $\tfrac14$ and $\tfrac18$, each gap half the one before. This suggests the formula
\[
x_n=1-2^{-n},
\]
where $2^{-n}=1/2^n$ as in Section~\ref{ch01:sec:algebra}. We prove it by induction on $n$.

For the base case, $n=0$, the formula gives $1-2^{0}=1-1=0$, which is $x_0$. For the induction step, let $n$ be a natural number and assume that $x_n=1-2^{-n}$. (The letter $n$ now plays the role of $k$; the name does not matter.) Substitute the induction hypothesis into the rule:
\begin{align*}
 x_{n+1}
 &=\frac{x_n+1}{2} &&\text{the rule for the next term}\\
 &=\frac{(1-2^{-n})+1}{2} &&\text{induction hypothesis}\\
 &=\frac{2-2^{-n}}2\\
 &=1-\frac{2^{-n}}2\\
 &=1-2^{-(n+1)}.
\end{align*}
The last equality uses $2^{-n}/2=1/(2^n\cdot2)=2^{-(n+1)}$. This is the formula at $n+1$, so by induction it holds for every natural number $n$.

The gap between $x_n$ and $1$ is exactly $1-x_n=2^{-n}$, which halves at every step, so the terms come as close to $1$ as we like; for instance, $x_{10}=1-1/1024$ is within $0.001$ of $1$. Yet the gap is never zero, so no term ever equals $1$.

\subsection*{Tracking the error instead}
Why is $1$ the right number to compare with? If some term were exactly $1$, the next term would be $(1+1)/2=1$ again. A number that an update rule leaves unchanged is called a \emph{fixed point} of the rule, and by Section~\ref{ch04:sec:unique}, $1$ is the only real number $x$ with $x=(x+1)/2$, so it is the only fixed point of this rule.

So instead of tracking $x_n$ itself, track its \emph{error} $x_n-1$, the difference between the current term and the fixed point; its absolute value $\abs{x_n-1}$ is the distance from $x_n$ to $1$. Subtract $1$ from both sides of the rule:
\[
x_{n+1}-1=\frac{x_n+1}{2}-1=\frac{x_n+1-2}{2}=\frac{x_n-1}{2}.
\]
Each update exactly halves the error, and this calculation never used the starting value $x_0=0$. So for the same rule with \emph{any} real starting value $x_0$, induction gives
\[
x_n-1=2^{-n}(x_0-1)\qquad\text{for every } n\in\NN.
\]
At $n=0$ both sides equal $x_0-1$, and if the formula holds at $n$, halving both sides gives it at $n+1$, by the rule just found; Exercise~\ref{ex:5.5} asks you to write out this induction in full. With $x_0=0$ the formula becomes $x_n-1=-2^{-n}$, that is, $x_n=1-2^{-n}$ again.

This more general statement covers every starting value at once, yet its induction step is shorter than the one above, because it records only what each step does, namely halve the error, and leaves out the starting value $0$. So when an induction proof gets tangled, try proving a more general statement; its induction step can be simpler. The contraction theorem of Chapter~\ref{ch:contraction} uses the same idea.
